\section*{All-dimensional supercritical sign}
The shift $m=u+3$ makes every coefficient of the numerator of $R_m$ positive,
and every coefficient of the positive polynomial inside $-C_m$ positive.
Thus $R_m>0$ and $C_m<0$ for $m\ge3$.

Use
\[
0\le \mathcal H_m\le\frac{m-2}{90m-179}.
\]
Because $C_m<0$,
\[
N_m\ge R_m+C_m\frac{m-2}{90m-179}.
\]
After clearing the positive denominator, the lower-bound numerator is
\[
\begin{aligned}
L_m={}&2729945147827667886720m^5
-27755132420474170999952m^4\\
&+112813395868533457497683m^3
-229153280695458887386228m^2\\
&+232620996871721820873517m
-94412163900120968220300.
\end{aligned}
\]
Its coefficients after $m=u+3$ are
\[
\begin{gathered}
2729945147827667886720,
13194044796940847300848,
25446870127333515303059,\\
24475321329207325509911,
11736484608366875120374,
2243933770033305838488,
\end{gathered}
\]
all strictly positive.  Hence $N_m>0$.  The adjoint calculation gives
$\ell^Tr<0$, so
\[
\boxed{c_m<0\qquad(m\ge3).}
\]
Together with $\eta_m>0$, the bifurcation is supercritical for every member of
the family.
